% ============================================================================
% BLOCK-003 Derivation v7: Normalization Candidate Catalog
% ============================================================================
% GUARDRAIL: This is a research catalog within paper_gravity_block003/ only.
% Do NOT expand scope to edc_book_2/ or modify FROZEN main.tex.
% ============================================================================
\documentclass[11pt,a4paper]{article}

\usepackage{fontspec}
\usepackage{amsmath,amssymb,amsthm}
\usepackage{physics}
\usepackage{graphicx}
\usepackage{booktabs}
\usepackage{array}
\usepackage{longtable}
\usepackage{xcolor}
\usepackage[hidelinks]{hyperref}
\usepackage[margin=2.5cm]{geometry}
\usepackage{fancyhdr}
\usepackage{tcolorbox}

% Epistemic tags
\newcommand{\BL}{\textsc{[BL]}}
\newcommand{\Ident}{\textsc{[I]}}
\newcommand{\Post}{\textsc{[P]}}
\newcommand{\Deriv}{\textsc{[Dc]}}
\newcommand{\Open}{\textsc{[Open]}}

% Status box
\newtcolorbox{statusbox}[1][]{
  colback=yellow!10,
  colframe=orange!80!black,
  title=#1,
  fonttitle=\bfseries
}

% Header/footer
\pagestyle{fancy}
\fancyhf{}
\fancyhead[L]{\small BLOCK-003 / Derivation v7}
\fancyhead[R]{\small EDC Gravity Program}
\fancyfoot[C]{\thepage}

\title{\textbf{EDC BLOCK-003 (Derivation v7):}\\[0.3em]
Normalization Candidate Catalog\\[0.3em]
\large How to Fix the Constant $C$ in $\kappa_5^2 = C \sigma^{-3/4}$?}
\author{EDC Working Document\\[0.5em]
\small 2026-02-02 --- Research note / catalog; not a results paper}
\date{}

\begin{document}
\maketitle

\begin{abstract}
Derivations v3--v5 established that the 5D gravitational coupling satisfies
$\kappa_5^2 = C \sigma^{-3/4}$ with the dimensionless constant $C$ unfixed
by brane tension $\sigma$ alone (No-Go Lemma). This note catalogs 5 candidate
normalization principles that could fix $C$ without circularly inserting the
observed Newton constant $G_N^{\rm obs}$. For each candidate, we state the
required extra input, assess circularity risk, and evaluate viability.
\textbf{Outcome:} Catalog complete; one candidate (NC-1: graviton zero-mode
normalization) identified as best next attempt. BLOCK-003 remains \Open.
\end{abstract}

\tableofcontents
\newpage

% ============================================================================
\section{Context Recap}
% ============================================================================

From derivations v1--v5:

\begin{itemize}
\item \textbf{v1:} Standard brane-world reduction yields $G_N = \kappa_5^2 / (6\pi L)$
for compact extra dimension of size $L$.

\item \textbf{v2:} Testing $L = R_\xi$ (weak scale) is dimensionally consistent
but leaves $\kappa_5^2$ unknown.

\item \textbf{v3/v4:} Dimensional analysis uniquely gives
$\kappa_5^2 = C \sigma^{-3/4}$, but $C$ cannot be fixed by $\sigma$ alone
(No-Go Lemma: only 2 independent scales available).

\item \textbf{v5:} Two clean options documented: NP1 (postulate $C = 8\pi$) or
NP2 (calibrate to $G_N^{\rm obs}$).
\end{itemize}

\textbf{The question:} Can we find a \textit{non-circular} principle that
fixes $C$ from EDC-internal or standard physics inputs, without fitting to
$G_N^{\rm obs}$?

% ============================================================================
\section{Anti-Circularity Definition}
\label{sec:circularity}
% ============================================================================

\begin{statusbox}[Anti-Circularity Criterion]
A normalization principle is \textbf{circular} if it:
\begin{enumerate}
\item Uses $G_N^{\rm obs}$ (or $M_{\rm Pl}$) as input to fix $C$.
\item Derives a scale that is itself defined via $G_N^{\rm obs}$.
\item Tunes a free parameter to match $G_N^{\rm obs}$ predictions.
\end{enumerate}
A principle is \textbf{non-circular} if $C$ is determined before any
comparison to gravitational observations.
\end{statusbox}

Note: Using $G_N^{\rm obs}$ for \textit{verification} (after $C$ is fixed)
is allowed; using it as \textit{input} is not.

% ============================================================================
\section{Candidate Catalog}
\label{sec:catalog}
% ============================================================================

\begin{table}[h]
\centering
\small
\begin{tabular}{llp{3.5cm}p{2.8cm}ccp{2.5cm}}
\toprule
\textbf{ID} & \textbf{Name} & \textbf{Fixes $C$ via} & \textbf{Extra Input} &
\textbf{Circ.} & \textbf{Status} & \textbf{Notes} \\
\midrule
NC-1 & Zero-mode norm. & Graviton KK reduction & $L$, BC type & LOW & Viable &
Best candidate \\
NC-2 & Induced gravity & DGP-like $R^{(4)}$ term & $N$, $\Lambda_{\rm UV}$ &
MED & Underdetermined & Species count unknown \\
NC-3 & RS warp matching & $M_{\rm Pl}^2 \sim \int e^{-2A}$ & $A(\xi)$, $k$ &
MED & Missing input & EDC doesn't specify $A$ \\
NC-4 & Topological quant. & Junction/Z$_6$ discrete & Topology data & LOW &
Speculative & No mechanism yet \\
NC-5 & Bulk $\Lambda_5$ match & $\Lambda_5 \sim \kappa_5^2 \rho_P$ &
$\Lambda_5$ value & HIGH & Near-circular & $\Lambda_5$ often tuned \\
\bottomrule
\end{tabular}
\caption{Normalization candidate catalog. Circularity risk:
LOW = no $G_N^{\rm obs}$ dependence;
MED = indirect dependence possible;
HIGH = likely circular.}
\label{tab:catalog}
\end{table}

% ============================================================================
\section{Candidate Details}
% ============================================================================

\subsection{NC-1: Graviton Zero-Mode Normalization}

\textbf{Idea:} In Kaluza-Klein reduction, the 4D graviton emerges as the
zero-mode of the 5D metric fluctuation. The 4D Planck mass is determined by
the normalization integral over the extra dimension.

\textbf{Setup:} Consider the 5D Einstein-Hilbert action
\[
S_5 = \frac{1}{2\kappa_5^2} \int d^5X \sqrt{-g^{(5)}} R^{(5)}.
\]
For a factorizable ansatz $ds^2 = g_{\mu\nu}(x) dx^\mu dx^\nu + d\xi^2$
with $\xi \in [0, L]$:
\[
S_5 \to \frac{L}{2\kappa_5^2} \int d^4x \sqrt{-g^{(4)}} R^{(4)} + \ldots
\]
Comparing with the 4D action $S_4 = (1/16\pi G_N) \int d^4x \sqrt{-g} R$:
\begin{equation}
\frac{1}{16\pi G_N} = \frac{L}{2\kappa_5^2}
\quad \Rightarrow \quad
G_N = \frac{\kappa_5^2}{8\pi L}
\label{eq:nc1_gn}
\end{equation}

\textbf{How it fixes $C$:} If we require the zero-mode to be \textit{canonically
normalized} (unit kinetic term), then:
\begin{equation}
\kappa_5^2 = 8\pi G_N L
\label{eq:nc1_kappa}
\end{equation}
Combined with $\kappa_5^2 = C \sigma^{-3/4}$:
\begin{equation}
C = 8\pi G_N L \cdot \sigma^{3/4}
\label{eq:nc1_c}
\end{equation}

\textbf{Circularity check:}
\begin{itemize}
\item If $L$ is specified independently (e.g., $L = R_\xi$), and $\sigma$ is
EDC-internal, then $C$ is fixed \textit{without} using $G_N^{\rm obs}$ as input.
\item However, Eq.~\eqref{eq:nc1_c} contains $G_N$. The question is: can we
express $G_N$ in terms of EDC quantities before comparison to observation?
\end{itemize}

\textbf{Resolution:} Eq.~\eqref{eq:nc1_gn} is a \textit{derived relation},
not an input. If we specify $L$ and $\kappa_5^2$ (with $C$ from convention),
then $G_N$ is a \textit{prediction}. The canonical normalization condition
(unit coefficient for $h_{\mu\nu}$ kinetic term) is a standard convention,
not a fit.

\textbf{Required input:} Compactification length $L$ (e.g., $L = R_\xi$ \Post).

\textbf{Status:} \textbf{Viable} --- best candidate for P18-B attempt.

\subsection{NC-2: Induced Gravity (DGP-like)}

\textbf{Idea:} Even without a fundamental 5D Einstein term, quantum loops of
brane-localized matter induce an effective $R^{(4)}$ term on the brane:
\[
S_{\rm induced} = \frac{M_4^2}{2} \int d^4x \sqrt{-h} R^{(4)}
\]
with $M_4^2 \sim N \Lambda_{\rm UV}^2 / (16\pi^2)$ in 4D.

\textbf{How it fixes $C$:} The induced 4D Planck mass competes with the
bulk-derived contribution. At long distances, the dominant term sets $G_N$.

\textbf{Problem:} The induced mass depends on:
\begin{itemize}
\item $N$: number of species (unknown in EDC)
\item $\Lambda_{\rm UV}$: UV cutoff (not specified)
\item Regularization scheme (introduces $\mathcal{O}(1)$ factors)
\end{itemize}

\textbf{Circularity check:} LOW if $N$ and $\Lambda_{\rm UV}$ are EDC-specified;
MED if they require external input or tuning.

\textbf{Status:} \textbf{Underdetermined} --- cannot fix $C$ without specifying
$N$ and $\Lambda_{\rm UV}$.

\subsection{NC-3: Randall-Sundrum Warp Matching}

\textbf{Idea:} In warped geometries with metric
$ds^2 = e^{-2A(\xi)} \eta_{\mu\nu} dx^\mu dx^\nu + d\xi^2$,
the 4D Planck mass is:
\[
M_{\rm Pl}^2 = \frac{1}{\kappa_5^2} \int_0^L d\xi\, e^{-2A(\xi)}
\]

\textbf{How it fixes $C$:} If $A(\xi)$ is specified, the integral is computable
and relates $\kappa_5^2$ to $M_{\rm Pl}^2$.

\textbf{Problem:} EDC does not currently specify:
\begin{itemize}
\item The warp factor $A(\xi)$
\item The AdS curvature scale $k$ (in RS, $A = k|\xi|$)
\end{itemize}

\textbf{Circularity check:} MED --- if $k$ is fit to reproduce correct $G_N$,
this becomes circular.

\textbf{Status:} \textbf{Missing input} --- requires EDC to specify bulk
geometry or warp factor.

\subsection{NC-4: Topological Quantization}

\textbf{Idea:} EDC's junction/Z$_6$ structure might provide discrete
normalization. For example, if the junction action has a topological
prefactor (winding number, Euler characteristic), this could quantize $C$.

\textbf{How it fixes $C$:} Discrete topology $\Rightarrow$ $C \in \{c_1, c_2, \ldots\}$
where $c_i$ are determined by topological data.

\textbf{Problem:} No concrete mechanism exists in current EDC formulation.
The Z$_6$ sector describes particle content, not gravitational normalization.

\textbf{Circularity check:} LOW (topology is internal).

\textbf{Status:} \textbf{Speculative} --- attractive in principle but no
derivation pathway available.

\subsection{NC-5: Bulk Cosmological Constant Matching}

\textbf{Idea:} The bulk cosmological constant $\Lambda_5$ is related to
$\kappa_5^2$ via Einstein equations. If $\Lambda_5$ is fixed by EDC
(e.g., from Plenum energy density), this could constrain $\kappa_5^2$.

\textbf{How it fixes $C$:} From $\Lambda_5 = -6k^2$ (AdS) and fine-tuning
relations, one might extract $\kappa_5^2$.

\textbf{Problem:} In standard RS, the relation is:
\[
\sigma = \frac{6k}{\kappa_5^2}
\]
which gives $\kappa_5^2 \sigma = 6k$, i.e., $C \sigma^{1/4} = 6k$.
This requires knowing $k$, which is often tuned to hierarchy requirements.

\textbf{Circularity check:} HIGH --- $\Lambda_5$ or $k$ is typically tuned
to reproduce desired $G_N$ or hierarchy.

\textbf{Status:} \textbf{Near-circular} --- unlikely to provide non-circular
normalization.

% ============================================================================
\section{Summary and Recommendation}
% ============================================================================

\begin{table}[h]
\centering
\begin{tabular}{llll}
\toprule
\textbf{ID} & \textbf{Viability} & \textbf{Missing Element} & \textbf{Next Step} \\
\midrule
NC-1 & Viable & $L$ specification & \textbf{Attempt in P18-B} \\
NC-2 & Underdetermined & $N$, $\Lambda_{\rm UV}$ & Requires EDC spectrum \\
NC-3 & Missing input & $A(\xi)$, $k$ & Requires bulk geometry \\
NC-4 & Speculative & Derivation mechanism & Future research \\
NC-5 & Near-circular & Non-tuned $\Lambda_5$ & Unlikely path \\
\bottomrule
\end{tabular}
\caption{Summary of candidate viability.}
\label{tab:summary}
\end{table}

\begin{statusbox}[Recommendation]
\textbf{Best next attempt: NC-1 (Graviton zero-mode normalization).}

Rationale:
\begin{itemize}
\item Lowest circularity risk (standard KK procedure, no $G_N$ input).
\item Only requires specifying $L$ (can use $L = R_\xi$ as \Post).
\item Produces a concrete expression: $C = 8\pi L \sigma^{3/4} G_N$, but
with $G_N$ as \textit{output}, not input.
\item Clear success/fail criterion: if $G_N^{\rm predicted}$ matches
$G_N^{\rm obs}$, the model is consistent; if not, it is falsified.
\end{itemize}

\textbf{P18-B task:} Derive $G_N$ from NC-1 with $L = R_\xi$, compute
numerical prediction, compare to observation.
\end{statusbox}

\textbf{BLOCK-003 Status:} Remains \Open. This catalog identifies the path;
execution is deferred to P18-B.

\end{document}
